\documentclass[12pt,a4paper,oneside]{report}
\usepackage[utf8]{inputenc}
\usepackage[T1]{fontenc}
\usepackage{lmodern}
\usepackage[margin=1in]{geometry}
\usepackage{setspace}
\usepackage{graphicx}
\usepackage{float}
\usepackage{caption}
\usepackage{booktabs}
\usepackage{tabularx}
\usepackage{array}
\usepackage{enumitem}
\usepackage{fancyhdr}
\usepackage[hidelinks]{hyperref}
\usepackage{url}

\onehalfspacing
\graphicspath{{figures/}}
\setlength{\parskip}{6pt}
\setlength{\parindent}{0pt}
\sloppy

\newcolumntype{Y}{>{\raggedright\arraybackslash}X}

\pagestyle{fancy}
\fancyhf{}
\fancyhead[L]{\small Brain Tumor Detection}
\fancyhead[R]{\small \thepage}
\renewcommand{\headrulewidth}{0.4pt}

\captionsetup{font=small,labelfont=bf}

\begin{document}


% --------------------------------------------------
% PAGE 1 : TITLE PAGE
% --------------------------------------------------

\begin{titlepage}
\centering

\vspace*{1cm}

{\Huge \textbf{Brain Tumor Detection from MRI Images using an Ensemble of CNN, Transfer Learning and Swin Transformer Models}}\\[1cm]

{\Large 7th/8th Semester Project Report}\\[1cm]

submitted to\\[0.3cm]

{\Large \textbf{Techno India University, West Bengal}}\\[1cm]

In partial fulfilment of the requirements\\
for the award of the degree of\\[0.3cm]

{\Large \textbf{Bachelor of Technology}}\\[0.3cm]

in\\[0.3cm]

{\Large \textbf{Computer Science \& Engineering}}\\[1cm]

by\\[0.5cm]

\begin{tabular}{c}
\textbf{Ronit Singh (University Roll No. - 221001001297)} \\[0.3cm]
\textbf{Ahana Sarkar (University Roll No. - 221001001222)} \\[0.3cm]
\textbf{Swikriti Roy (University Roll No. - 221001001279)} \\[0.3cm]
\textbf{Dibyashree Dey (University Roll No. - 221001001346)} \\[0.3cm]
\textbf{Shreyansh Kumar Gupta (University Roll No. - 221001001290)}
\end{tabular}

\vspace{1.5cm}

Under the guidance of\\[0.3cm]

{\Large \textbf{Dr. Priyanka Saha}}\\[1.5cm]

Department of Computer Science \& Engineering\\
\textbf{TECHNO INDIA UNIVERSITY, WEST BENGAL}\\[1cm]

January 2026\\[0.5cm]

\copyright\ 2026, Techno India University. All rights reserved.

\end{titlepage}

% --------------------------------------------------
% PAGE 2 : BLANK PAGE
% --------------------------------------------------

\newpage
\thispagestyle{empty}
\mbox{}

% --------------------------------------------------
% PAGE 3 : DEDICATION
% --------------------------------------------------

\newpage
\thispagestyle{empty}

\vspace*{7cm}

\begin{center}
{\Huge Dedicated to}\\[1cm]
Our parents, families, and teachers who supported us throughout this journey.
\end{center}

% --------------------------------------------------
% PAGE 4 : BLANK PAGE
% --------------------------------------------------

\newpage
\thispagestyle{empty}
\mbox{}

% --------------------------------------------------
% PAGE 5 : DECLARATION OF AUTHORSHIP
% --------------------------------------------------

\newpage

\begin{center}
{\Large \textbf{DECLARATION OF AUTHORSHIP}}
\end{center}

\vspace{1cm}

We hereby declare that the project report entitled ``Brain Tumor Detection from MRI Images using an Ensemble of CNN, Transfer Learning and Swin Transformer Models'' is an authentic
record of our own work carried out at the Department of Computer Science and
Engineering, Techno India University, West Bengal, during the 8th semester
of the academic year 2025--2026 under the supervision of Dr. Priyanka Saha,
Assistant Professor.

\vspace{0.5cm}

We further declare that the matter embodied in this project report has not been submitted
by us for the award of any other degree or diploma of this or any other Institute/University.
All the information have been obtained and presented in accordance with academic
rules and ethical conduct. We also declare that, as required by these rules and conduct,
we have fully cited and referenced all materials and results that are not original to this
work. The dataset used in this project is publicly available under its original license; no
patient-identifiable data was collected, generated, or stored.

\vspace{1cm}

\textbf{Place:} EM-4/1, Sector V, Bidhannagar, Kolkata, West Bengal -- 700091

\vspace{0.5cm}

\textbf{Date:} January 2026

\vspace{1cm}

\textbf{Signatures:}

\vspace{1cm}

\begin{tabular}{p{7cm}p{7cm}}
Signature: & Signature: \\
Name: Ronit Singh & Name: Ahana Sarkar \\
Roll No: 221001001297 & Roll No: 221001001222
\end{tabular}

\vspace{1cm}

\begin{tabular}{p{7cm}p{7cm}}
Signature: & Signature: \\
Name: Swikriti Roy & Name: Dibyashree Dey \\
Roll No: 221001001279 & Roll No: 221001001346
\end{tabular}

\vspace{1cm}

\begin{tabular}{p{7cm}}
Signature: \\
Name: Shreyansh Kumar Gupta \\
Roll No: 221001001290
\end{tabular}

% --------------------------------------------------
% PAGE 6 : BLANK PAGE
% --------------------------------------------------

\newpage
\thispagestyle{empty}
\mbox{}

% --------------------------------------------------
% PAGE 7 : CERTIFICATE OF RECOMMENDATION
% --------------------------------------------------

\newpage

\begin{center}
{\Large \textbf{CERTIFICATE OF RECOMMENDATION}}
\end{center}

\vspace{1cm}

This is to certify that the work embodied in this thesis entitled ``Brain Tumor Detection from MRI Images using an Ensemble of CNN, Transfer Learning and Swin Transformer Models''
has been satisfactorily completed by Ronit Singh (Roll No. 221001001297),
Ahana Sarkar (Roll No. 221001001222),
Swikriti Roy (Roll No. 221001001279),
Dibyashree Dey (Roll No. 221001001346),
and Shreyansh Kumar Gupta (Roll No. 221001001290).

It is a bonafide piece of work
carried out under my supervision and guidance at Techno India University, Kolkata,
for partial fulfilment of the requirements for the awarding of the Bachelor of Technology
in Computer Science \& Engineering degree of the Department of
Computer Science and Engineering, Techno India University, during the academic year
2025--2026.

\vspace{2cm}

\textbf{Dr. Priyanka Saha}\\
Assistant Professor, Department of Computer Science and Engineering,\\
Techno India University, Kolkata, West Bengal, India.\\
(Supervisor)

\vspace{2cm}

\textbf{Forwarded By:}

\vspace{1cm}

\textbf{[Name of Head of Department]}\\
HoD, Department of Computer Science and Engineering,\\
Techno India University, Kolkata, West Bengal, India.

% --------------------------------------------------
% PAGE 8 : BLANK PAGE
% --------------------------------------------------

\newpage
\thispagestyle{empty}
\mbox{}

% --------------------------------------------------
% PAGE 9 : ACKNOWLEDGEMENT
% --------------------------------------------------

\newpage

\begin{center}
{\Large \textbf{ACKNOWLEDGEMENT}}
\end{center}

\vspace{1cm}

We would like to first express our sincere gratitude to our project supervisor,
Dr. Priyanka Saha, Assistant Professor, Department of Computer Science and Engineering,
Techno India University, West Bengal. Their continued guidance, technical insight, and
constructive feedback during every stage of this project were instrumental in shaping the
final outcome. Their willingness to discuss difficult engineering decisions, including the
diagnosis of subtle numerical-stability issues during mixed-precision training, helped us
approach the work with both rigour and curiosity.

\vspace{0.5cm}

We take this opportunity to express our gratitude to all faculty members of the Department
of Computer Science and Engineering for their support, the encouragement they provided
during our coursework, and the foundation they laid in machine learning, computer vision,
and software engineering.

\vspace{0.5cm}

We acknowledge the open-source community whose tools made this work feasible: the PyTorch
and timm library maintainers, the authors of pytorch-grad-cam, the Streamlit and FastAPI
teams, and Masoud Nickparvar for releasing the Brain Tumor MRI Dataset on Kaggle under a
permissive license that supports academic research.

\vspace{0.5cm}

Finally, we thank our parents and families for their unceasing encouragement, patience,
and support throughout the semester.

\newpage
\thispagestyle{empty}
\mbox{}

\pagenumbering{roman}

\chapter*{Abstract}
\addcontentsline{toc}{chapter}{Abstract}
Tumor classification in MRI scans of the brain is an essential application in medical science but requires time, expertise, and manual interpretation, leading to discrepancies between observers. It has been reported that the accuracy of observer-based tumor types varies within the range of 70 to 85 percent, and expert neuroradiologists' accessibility differs according to geographic location. Thus, there is a need for an AI-powered assistance solution that provides fast, consistent, and explainable tumor classifications.
The current project proposes an end-to-end deep learning model that classifies brain MRI scans (T1) into four tumor categories: glioma, meningioma, pituitary tumor, and no tumor. An ensemble of three different architectures is proposed as the classification technique. The first model is a custom-designed CNN with about 1.2 million parameters that undergoes training from a random initialization process. EfficientNet-B0 is another CNN with around 4.0 million parameters which is fine-tuned using two-step transfer learning. Finally, Swin-Tiny is a vision transformer with 27.5 million parameters that applies a window-based self-attention scheme in its architecture. All three models have their softmax layer output concatenated using a weighted average model where mixing weights are tuned by grid search on the validation dataset.
It was trained and validated on the open-source Brain Tumor MRI Dataset (Nickparvar, Kaggle, 7,200 images belonging to four categories). An 80/20 stratified split divides the 5,600 training images into 4,480 training data and 1,120 validation examples, leaving 1,600 testing images held out and evaluated just once. In this held-out test set, the best single model (EfficientNet-B0) obtains an accuracy score of 94.94\% and a macro-average ROC-AUC score of 0.991, whereas the best three-model ensemble produces 94.69\% accuracy at 0.991 AUC. Here is one straightforward empirical observation: ensembling did not improve upon but matched the performance of the dominant base learner. This phenomenon will be explained using the theory of dominant-model ensembles.
In addition to accuracy, our system generates per-prediction Grad-CAM maps for each of the three models, a low-confidence thresholding gate which makes an explicit recommendation for radiologist evaluation of the most ambiguous cases, and four deployment options: Streamlit, FastAPI, ONNX inference with a 2-3.5X CPU runtime speed-up, and Docker containerization.
Keywords: Brain Tumor Classification, Magnetic Resonance Imaging, Convolutional Neural Network, Transfer Learning, Swin Transformer, Ensemble Learning, Grad-CAM, Computer-Aided Diagnosis.
\tableofcontents
\listoffigures
\listoftables
\clearpage
\pagenumbering{arabic}
\chapter{Chapter 1: Introduction}
\section{1.1 Medical Imaging and Computer-Aided Diagnosis}
\begin{figure}[H]
  \centering
  \includegraphics[width=0.92\textwidth]{Figure_08_system_flow_diagram.png}
  \caption{High-level system flow: MRI input, preprocessing, three parallel models, weighted ensemble, predicted class with confidence and Grad-CAM heatmaps.}
\end{figure}
Brain tumours refer to unusual cell proliferation within the brain. The cells could be from brain tissues themselves (primary tumours) or could have spread to the brain from other body organs (metastatic tumours). According to the World Health Organization (WHO), there are over 150 different types of brain tumours depending on the cell type involved, molecular characteristics, and the level of cancer. Of the over 150 types, gliomas, meningiomas, and pituitary adenomas account for the majority of the cases seen clinically.
Magnetic resonance imaging (MRI) is considered the most effective diagnostic method due to its high-quality image resolution and no exposure to ionising radiation. On average, a normal brain MRI scan yields numerous two-dimensional images through several sequences (T1-weighted, T2-weighted, FLAIR, and contrast-enhanced T1) that must all be examined by a neuroradiologist. On average, a neuroradiologist takes between 15 and 30 minutes per patient when analysing brain images manually, with an observer accuracy rate of about 70 to 85 per cent [1].
Computer-Aided Diagnosis (CAD) seeks to be helpful to radiologists as they obtain computerized second reader interpretations. Deep learning technology and specifically Convolutional Neural Networks (CNN) have become a game changer for the CAD industry over the last decade. Well-tuned models utilizing such an algorithm can process an MRI slice in milliseconds, give a highly accurate prediction and even point out on which pixels in the image the model relied the most while making its decision. As such, a computer-aided diagnosis based on deep learning could be used effectively as a triaging aid in high-traffic radiology clinics and as a decision-making aid in regions where access to specialist radiologists is limited [2].
According to the International Association of Cancer Registries and the Indian Journal of Neurology (2025), there are about 28,000 brain cancer patients diagnosed in India each year, though some sources estimate that the number may reach up to 40-50,000 annually. Brain and central nervous system tumors account for 2\% to 3\% of all types of cancers in the country. At the same time, the number of radiologists in India is less than one per 100,000 people. The deficit is especially noticeable in rural and semi-urban territories [17].
\section{1.2 Objectives of the Project}
\subsection{1.2.1 Primary Objective}
Creation of a brain tumor classifier using deep-learning methods, which receives as input a T1-weighted brain MRI slice and outputs a prediction about the presence of glioma, meningioma, pituitary tumors, or absence of any tumor with high accuracy, well-calibrated confidence, and visual explanations, using three different architectures of neural networks in combination.
\subsection{1.2.2 Secondary Objectives}
Objective 1: Creation of a baseline CNN classifier from scratch for benchmarking and providing diversity in neural network architectures in the ensemble.
Objective 2: Fine-tuning of an EfficientNet-B0 pre-trained network on the task of brain tumor classification using a two-stage transfer learning approach (freezing-then-fine-tuning), which achieves maximum accuracy without the risk of catastrophic forgetting of pretrained weights.
Objective 3: Fine-tuning of a Swin-Tiny vision transformer on the same task to bring architectural diversity in neural network structures through self-attention mechanism.
Objective 4: Combine the three trained models into an ensemble using weighted averages where the weights are calculated by exhaustively searching through a grid on the validation set, and compare it to the performance of the soft voting and stacking baseline ensembles.
Objective 5: Implement Gradient-weighted Class Activation Mapping (Grad-CAM) for all three of our models so we can explain the regions that our models focused on when making their predictions in visual form.
Objective 6: Develop a full-scale production application that features a Streamlit GUI frontend with a FastAPI REST backend service to enable inference on our models in real time, along with Grad-CAM output.
Objective 7: Export all three of our models to the ONNX format to make them deployable across platforms, and evaluate their speed and ability to perform accurate classification on the CPU.
Objective 8: Perform a thorough evaluation of the entire pipeline on 1,600 image holdout test set that is never used for training or hyperparameter tuning and report on accuracy metrics like F1-Score, macro-ROCAUC, per-class confusions matrices, and inference times.
\section{1.3 Challenges in Brain Tumor Classification}
Automated brain tumor classification using MRIs poses some difficulties that any effective solution will have to contend with:
\begin{itemize}[leftmargin=*]
  \item Intra-class variation: The same class of tumor can look extremely different between patients based on its size, location, development phase, and scan parameters. A glioma tumor will look completely different from another in another patient.
  \item Visual similarity between classes: Some tumors are exceptionally similar in appearance to others when only one T1 sequence is considered. Gliomas and meningiomas, for example, could look almost identical to each other at first glance even for expert radiologists. Further imaging is needed to differentiate between the two classes.
  \item Few training images: The number of medical imaging data used to train models is always limited because it is highly sensitive information. Unlike the millions of natural images available for ImageNet, for example, only about 5,600 training images are available here.
  \item Heterogeneity of scanners and protocols: Images obtained using various scanners, with varying magnetic field strength, coils, and imaging protocols lead to systematically different intensity distributions. An algorithm that has been developed using images taken with one specific machine will most likely exhibit inferior performance when dealing with images collected with a different one.
  \item Explainability and confidence in diagnosis: Radiologists will be unlikely to accept a diagnostic aid as reliable if it is completely unexplainable. If the brain tumor detection system is expected to become clinically applicable, then it should be able to deliver its conclusions along with explanations as to why a certain result was predicted. In this respect, Grad-CAM maps satisfy the requirement.
  \item Uncertainty calibration: An uncalibrated neural network that produces very confident incorrect results is more problematic than one producing uncertain predictions.
\end{itemize}
\section{1.4 Methodologies}
The methodology adopted in this research project comprises a series of stages ranging from data collection and preprocessing to model training, ensemble creation, interpretability, and application. These are presented in the list below, with full descriptions given in the respective later chapters.
\begin{itemize}[leftmargin=*]
  \item Data Collection: The Brain Tumour MRI Data Set (Nickparvar, Kaggle) is adopted, which includes 7,200 labelled MRI images divided into four categories. The data set is randomly divided into training (4,480 images), validation (1,120 images) and test sets (1,600 images) using stratification and a predefined seed value for consistency.More information can be found in Chapter 3.
  \item Data Preprocessing and Augmentation: All data is preprocessed to convert to the RGB format, rescaling to 224 by 224 pixels, and normalised according to ImageNet parameters. Training data is further augmented through random image flipping, random rotation, affine transformation, brightness/contrast adjustments, and random erasing. Refer to Chapter 3.
  \item Training Models: Three architecturally different models undergo training - a custom-made 4-block CNN (1.2 million parameters), an EfficientNet-B0 with two-stage transfer learning (4.0 million parameters), and a Swin Tiny transformer with two-stage fine tuning (27.5 million parameters). All three share identical AdamW optimizer, CosineAnnealingLR scheduler, bfloat16 mixed precision, and early stopping criteria. Further details can be found in Chapter 5.
  \item Building Ensemble: After training, each model outputs a 4-class softmax probability vector. Three ensemble combinations are considered - soft voting, averaging, and stacking. Weights for the averaging approach are estimated through an exhaustive search over the simplex at 0.05 step size (231 iterations). Further details can be found in Chapter 5.
  \item Explainability: For each model, we create grad-cam maps to visually identify what region of the MRI is more likely responsible for the predicted class. Both overlay images are integrated into the Streamlit UI and FastAPI interface. Further details can be found in Chapter 5.
  \item Deployment: The trained models are deployed within four wrappers - Streamlit web application for interactive usage, FastAPI REST API for programmatic access, ONNX export for interoperability purposes, and a Docker image for deployment to the cloud. Details can be found in Chapters 4 and 5.
\end{itemize}
\chapter{Chapter 2: Literature Survey}
\section{2.1 Introduction to the Field}
Recent systematic reviews offer an appropriate introduction. Bouhafra and El Bahi conducted a review of 60 studies dealing with brain tumor classification through deep learning during 2020-2024 using transfer learning, autoencoder, transformer, and attention techniques [1]. A comprehensive survey was conducted by Ranjbarzadeh et al. about Explainable Artificial Intelligence (XAI) and vision-transformer models applied to brain tumors [2]. The above works can help us distinguish three main research streams: (1) transfer learning from pre-trained ImageNet backbone models (EfficientNet and ResNet models); (2) vision transformer models (Swin series); and (3) architecture ensembling.
In section 2.2, we will introduce in detail the methods, data used, main findings, and contributions of 10 papers that have a direct bearing on our own project. In section 2.3, all these reviewed works will be brought together into one comparison table.
\section{2.2 Review of Related Work}
\subsection{2.2.1 Babu Vimala et al. (2023): Hybrid Deep Learning with EfficientNet Variants}
For the three-class brain tumor classification (glioma, meningioma, pituitary) of images obtained via CE-MRI Figshare dataset, Babu Vimala et al. used a technique based on fine-tuned versions of EfficientNet models (B0 through B4). The methodology involved a two-fold process by initializing EfficientNet models with weights from the ImageNet database and further appending custom layers to classify tumors. Visualization of attention maps using Grad-CAM was used. The most efficient of all models was EfficientNetB2 with 99.06\% test accuracy and an F1-score of 98.79\%. The research proved that the application of EfficientNet models in classification of medical images is feasible when combined with proper fine tuning and data augmentation [3].
\subsection{2.2.2 Islam et al. (2024): BrainNet Optimised EfficientNet Architecture.}
Islam et al. introduced BrainNet, an efficient architecture for the EfficientNet pipeline which was validated using the CE-MRI dataset consisting of 3,064 images of T1-weighted images. In their study, they evaluated all versions of EfficientNet from B0 to B7 and used sophisticated data pre-processing and data augmentation. The most effective architecture of the EfficientNet pipeline is EfficientNetB3 which was able to achieve an accuracy of 99.69\% in testing. According to their research findings, learning-rate scheduling and progressive unfreezing were key strategies to obtaining the best performance in small medical imaging datasets.
\subsection{2.2.3 Alnowami et al. (2024): Attention Mechanisms for EfficientNetv2}
Alnowami et al. incorporated attention modules Global Attention Mechanism (GAM) and Efficient Channel Attention (ECA) in EfficientNetv2 for classifying brain tumours from MRIs. With this approach, EfficientNetv2 was capable of detecting salient features within intricate MRIs to accurately classify them into one of the four categories with 99.76\% accuracy. The paper further introduces the use of Grad-CAM to ensure clinical interpretability in MRI analysis. It shows that CNN models supplemented with attention mechanisms can match fully-fledged transformers in performance but with significantly lower resource consumption, making the choice of EfficientNet-B0 over heavy transformer models reasonable under the constraint of 4 GB VRAM [5].
\subsection{2.2.4 Liu et al. (2021): Original Swin Transformer}
The authors presented the Swin Transformer, a vision transformer architecture that leverages shifted windows for self-attention computations. The novelty of this vision transformer lies in its computation approach, where instead of performing self-attention operations on the entire input image (a quadratic operation), the input image is segmented into disjoint windows, and attention is performed inside each window. Every other layer performs shifting of the windows by half of their width, enabling information flow between different windows. As a result, the Swin Transformer enjoys a linear computational complexity relative to the image size while maintaining global modelling capabilities. The tiny version of the Swin Transformer reaches 81.3\% accuracy in ImageNet with just 28 million parameters.
\subsection{2.2.5 Haq et al. (2024): Swin Transformer with HSW-MSA and ResMLP}
In their work, Haq et al. proposed a modified version of the Swin Transformer for brain tumour diagnosis . They used Hybrid Shifted Windows Multi-Head Self-Attention (HSW-MSA) and Residual MLP (ResMLP) instead of the traditional MLP layer. Modifications were made to increase accuracy, reduce memory requirements and reduce the complexity of training. The proposed architecture was examined on the open-source four-class brain MRI dataset by applying transfer learning and data augmentation techniques. As a result, the authors managed to reach 99.92\% accuracy, significantly outperforming the state-of-the-art CNN architectures. This study showed the great promise of transformers in medical imaging and thus inspired us to use Swin-Tiny as one of three base models in our ensemble [6].
\subsection{2.2.6 Alsubai et al. (2024): Hybrid Swin Transformer and ResNet50V2}
Authors proposed a new model based on the integration of Swin Transformer with ResNet50V2 for better classification of brain tumours using MRIs.The Swin Transformer model learns local and global hierarchical representations via sliding windows in self-attention mechanisms, whereas ResNet50V2 learns additional deep convolutional representations. The fusion is achieved using transfer learning where pretrained weights from both models are applied. The hybrid model proves that integrating both CNN-based local representations and transformer-based global context models produces excellent results under various MRI imaging conditions.
\subsection{2.2.7 Khan et al. (2023): Deep Learning Weighted Average Ensemble}
The authors proposed an ensemble technique containing a weighted average model. There are three deep learning feature spaces, VGG19 model (using transfer learning), traditional CNN model and CNN with data augmentation. The best combination of weights was found using an exhaustive grid search algorithm. The dataset employed was the TCGA lower-grade gliomas database, which contained 3,929 MRI brain images. The authors showed that the performance of the ensemble model is better than that of each individual component model regarding accuracy, precision, and F1-score. This study provides direct motivation for our project in terms of ensemble weights tuning [8].
\subsection{2.2.8 Abdella et al. (2025): Ensemble of seven CNNs with Majority Voting}
Abdella et al. compared seven pretrained architectures of convolutional neural networks, such as GoogLeNet and Inception-v3, used for brain tumor classification into four classes based on the T1-weighted magnetic resonance images. Both networks were trained utilizing two optimization algorithms, namely SGDM and Adam, while evaluation was conducted using a publicly available dataset divided into training, validation, and test subsets accounting for 70\%, 10\%, and 20\% of images correspondingly. Majority voting ensemble of all 14 models resulted in 99.8\% overall accuracy and AUC higher than 0.997 for all classes of tumors, outperforming all individual models. The diversity of models is the most important factor in ensembles, which is worth noting and is consistent with the strategy of our project [9].
\subsection{2.2.9 Ali et al. (2024) Clinical Explainability of Grad-CAM with ResNet50}
Ali et al. [27] proposed the use of Grad-CAM with ResNet50 to solve the major problem of explainability in brain tumour detection using deep learning. Ali et al. proposed using Grad-CAM with ResNet50. Their method trained the ResNet50 model with data augmentation on an MRI dataset, then used Grad-CAM to create a heatmap of parts of the MRI that were most significant to the model's prediction. This resulted in a model with 98.52\% accuracy and precision-recall performance greater than 98\%. This study shows that the integration of high accuracy together with explanations contributes to higher clinical confidence and helps radiologists validate that the model pays attention to appropriate disease markers. Our study uses the same technique of Grad-CAM for all three base models and the target class of the ensemble [10].
\subsection{2.2.10 Naser and Deen (2023): Multi-CAM Comparative Study}
A study conducted by Naser and Deen proposed a Lightweight Explainable Deep Learning model using three explainability methods, including CAM, Grad-CAM and Grad-CAM++ across three pre-trained deep learning architectures (pretrained VGG-19, scratch VGG-19 and EfficientNet). The results from experiments conducted on two publicly available benchmark MRI datasets, both multi-class and binary, indicate that the combination of pretrained VGG-19 architecture with Grad-CAM resulted in better accuracy and localisation heatmap performance than any other combination. This paper is among the limited number of works that conduct quantitative comparison of CAMs in medical applications, indicating that Grad-CAM (utilised in our project) provides clinically valuable heatmaps [11].
\section{2.3 Summary of Reviewed Approaches}
The Table below includes all ten reviewed papers in one place, highlighting the method used, datasets, overall accuracy, and the particular contribution to our project idea in the last column.
Table 2.1: Comparison summary of the literature being reviewed.
\begin{center}
\footnotesize
\begin{tabularx}{\textwidth}{YYYYY}
\toprule
\textbf{Reference} & \textbf{Year} & \textbf{Method} & \textbf{Accuracy} & \textbf{Contribution to our project} \\
\midrule
Babu Vimala et al. & 2023 & EfficientNetB0-B4 + Grad-CAM & 99.06\% & Two-step transfer learning baseline \\
Islam et al. & 2024 & EfficientNetB0-B7 with augmentation & 99.69\% & Progressive unfreezing schedule \\
Alnowami et al. & 2024 & EfficientNetv2 + GAM + ECA & 99.76\% & Attention can rival transformers \\
Liu et al. & 2021 & Swin Transformer (foundational) & 81.3\% (ImageNet) & Shifted-window self-attention \\
Haq et al. & 2024 & Modified Swin + HSW-MSA + ResMLP & 99.92\% & Transformer-on-MRI feasibility \\
Alsubai et al. & 2024 & Swin + ResNet50V2 hybrid & high & CNN + transformer fusion idea \\
Khan et al. & 2023 & Weighted ensemble VGG19 + CNN & best of 3 & Grid-search ensemble weights \\
Abdella et al. & 2025 & Majority voting (7 CNNs) & 99.8\% & Diversity drives ensembles \\
Ali et al. & 2024 & ResNet50 + Grad-CAM & 98.52\% & Clinical-trust framework \\
Naser \& Deen & 2023 & VGG-19 + CAM / Grad-CAM / Grad-CAM++ & high & Grad-CAM beats CAM and Grad-CAM++ \\
\bottomrule
\end{tabularx}
\end{center}
\section{2.4 Market Research}
The prevalence of brain tumors around the world is significant and increasing. Based on data provided by the Global Cancer Observatory (GLOBOCAN 2022), about 308,000 patients have been diagnosed with brain tumors or tumors of central nervous systems annually in 2020. The incidence rate for these types of diseases per 100,000 people worldwide is estimated at 5.57\%. Tumors of the brain and central nervous system occupy the 19th place in terms of popularity and the 12th place concerning mortality rates [17].
The situation with the burden of such diseases in India is even more critical. Thus, the International Association of Cancer Registries and the Indian Journal of Neurology estimate that annually, between 28,000 and 40,000-50,000 new brain tumors cases occur. The most affected population segment are children between 0 and 14 years old, where brain tumors are only secondary to lymphoid leukaemia. Tumors of the central nervous system in India have an incidence rate of 5 to 10 per 100,000 people [17].
With this huge disease burden, India lacks diagnostic specialists by the dozen. Less than one radiologist for every 100,000 Indians, with the gap particularly pronounced in tier-2, tier-3 Indian cities, and rural areas where costly imaging equipment remains underutilized as there is nobody with requisite skills to interpret these images. In such conditions, AI-powered diagnostic systems, capable of generating a reliable second opinion instantly, are a practical necessity [17].
The medical image analysis software market, fueled by advancements in the field of deep learning and digitization of radiology operations, is estimated to witness considerable growth until 2030. Governments across countries in the Asia-Pacific region, especially in India (Ayushman Bharat initiative) and China (Healthy China 2030 strategy), see diagnostic imaging and AI assistance in their analysis as an important health policy goal.
\section{2.5 Research Gap and Limitations of Existing Systems}
Despite the large body of work, there are several limitations in the literature:
Overstated accuracy: Many studies report only validation-set accuracy without a truly held-out test set, so headline numbers are inflated. This decrease in accuracy from val to test, which we find to be between 4 and 8 percentage points across our models, is rarely discussed.
Limited architectural diversity in the ensembles: Most of the ensemble studies combine variants of the same architecture family (e.g. multiple ResNets or multiple EfficientNets). Very few works intentionally combine architecturally different model families (CNN, depthwise-separable CNN, and transformer) to maximise the error diversity.
Lack of explainability: Grad-CAM is cited in several papers, but few compare heatmaps per model or analyse failure modes in the context of explainability. 
No deployment engineering:Most published work only goes up to model training and metric reporting. Few provide a working web interface, REST API or containerised deployment that a hospital IT team could realistically adopt.
No honest negative results:  The ensemble learning literature never reports ensembles doing worse than individual models. The case where a dominant base learner makes ensembling counterproductive is rarely acknowledged.
\section{2.6 Motivation for Project Selection}
This project was chosen based on the confluence of an urgent clinical need (reliable support for brain tumour classification), a well-characterised public dataset, and the desire to help fill the research gaps described above. The specific questions the project seeks to answer are:
Can combining 3 architecturally different models (custom CNN, transfer learning with depthwise-separable CNN, and vision transformer) outperform any single model in classification?
How big is the gap between validation-set and test-set accuracy?What does this gap say about the reliability of widely reported accuracy figures in the literature?
Are Grad-CAM heatmaps medically meaningful visual explanations? Do different model architectures focus on different anatomical regions?
Is it possible to wrap up the whole system from image upload to prediction with explainability as a deployable web application for demonstration in clinical settings?
Answering these questions, the project contributes not only a technical artifact (the trained models and deployment pipeline) but also empirical insights, in particular the documentation of an honest negative ensembling result, that have pedagogical and practical value.
\chapter{Chapter 3: Dataset and Data Preparation}
\section{3.1 Overview}
A thoroughly characterized and replicable dataset is the key to building an effective supervised deep learning model, and the validity of all accuracy statistics discussed in this report hinges upon the quality, representativeness, and origin of the dataset that was used for training and validation. This chapter discusses the dataset origin, the clinical interpretation of each class, the reasons for choosing this dataset, exploration, the dataset split into train/val/test sets, preprocessing, and augmentation.
\section{3.2 Dataset Source and Acquisition}
The dataset utilized for the entire research is the Brain Tumor MRI Dataset, which is an openly available data set provided by Masoud Nickparvar through Kaggle. The URL for this data set is:
\url{https://www.kaggle.com/datasets/masoudnickparvar/brain-tumor-mri-dataset}
The dataset is a compilation and curation of three other publicly available MRIs for the brain, which include figshare brain tumor dataset, SARTAJ brain tumor dataset, and Br35H brain tumor detection dataset, and these have been re-categorized by the researcher into a common taxonomic hierarchy based on four classes, with a common Train/Test file structure for each category [12].
The images have been made available in JPEG format in a two-level split/class/ directory structure. There are no DICOM headers, patient identification details, scanner details, and any other protected health information accompanying the images, which are nothing but 2D Axial Brain MRIs, anonymized at source, one per file. The balanced classes by the original author consist of 7,200 images in total out of which 5,600 images belong to the Training folder and 1,600 to the Testing folder.
\section{3.3 Class Categories and Clinical Meaning}
\begin{figure}[H]
  \centering
  \includegraphics[width=0.85\textwidth]{Figure_02_sample_mri_per_class.jpg}
  \caption{Representative MRI slices, one per class.}
\end{figure}
The four classes found in this dataset include the three most prevalent types of primary brain tumors, alongside the absence of any tumor.
Glioma: Gliomas are tumors that occur in the glial cells of the central nervous system, mainly in the astrocytes, oligodendrocytes, and ependymal cells. Gliomas represent the most prevalent malignant primary brain tumor in adult patients. The WHO classification system uses a grading system where gliomas can be graded from I, which is least aggressive, up to IV, which is highly aggressive. An example of a highly aggressive brain tumor is glioblastoma multiforme (grade IV), which is considered one of the most aggressive cancers in humans.
Meningioma:Meningiomas are tumors of the arachnoid cap cells of the meninges, which consist of three membranous coverings around the brain and spinal cord. The tumors account for the most common types of primary brain tumors, which tend to be benign and slow-growing. In MRI scans, meningiomas generally have a well-defined mass outside the gray matter or axioms and are associated with the characteristic appearance of the dural tail on contrast. Although meningiomas have a unique appearance on multimodal images, their appearance on a single axial plane of T1 can resemble that of a glioma.
Pituitary tumor:The term 'pituitary adenoma' refers to tumors developing from the pituitary gland, which is a small endocrine gland located at the base of the brain in a structure called the sella turcica. They are mostly non-cancerous tumors. When viewed using MRI, pituitary lesions can be found within or slightly above the sella turcica region, which is the base of the brain. Because of their restricted anatomic location, they are easy to classify using an automated classifier.
Absence of tumor: This category includes images that show no presence of a mass lesion. It is crucial for this group to be included in the study because an effective diagnostic tool should be able to accurately determine normal results since erroneous detection of tumors has its clinical consequences as well. 
\section{3.4 Rationale for Dataset Selection}
There exist several publicly available brain MRI databases, such as the BraTS challenge database, the figshare database, the IXI database, and other one-source databases. The Nickparvar database was chosen after analyzing each option based on the following factors:
A four-class taxonomy was consistent with the clinical issue. The majority of alternatives to the dataset only concern binary classification of tumors or non-tumors, or volumetric segmentation instead of classification on slices.
Class balance. After performing per-class augmentations according to the instructions of the original author, there is an almost equal amount of images in each class within the Training directory.
Sufficient volume for transfer learning. A total of roughly 7,200 images in the dataset make this volume sufficient for fine-tuning ImageNet pre-trained backbone networks without going into the very-low-data regime.
Public and liberal license. The dataset does not contain any identifiers of patients and can be used for academic purposes.
Baseline. Over the years since its creation in 2021, the database has featured in numerous peer-reviewed papers and, thus, offers us an effective benchmark to compete against.
Variability. Due to its aggregation of data from three other databases upstream, the images vary in their scanner type and intensities, making the model highly adaptive to inter-scanner variability.
\section{3.5 Dataset Statistics}
\begin{figure}[H]
  \centering
  \includegraphics[width=0.8\textwidth]{Figure_01_class_distribution.png}
  \caption{Per-class image counts across the Training, Validation and Test splits.}
\end{figure}
As seen from Table 3.1, the total dataset comprises the following numbers for each class. The Training directory is divided into training and validation datasets using stratified random sampling with a 80\% to 20\% split, respectively, while the Testing directory is kept for testing purposes only.
Table 3.1: Per-class image counts.
\begin{center}
\footnotesize
\begin{tabularx}{\textwidth}{YYYYYY}
\toprule
\textbf{Split} & \textbf{glioma} & \textbf{meningioma} & \textbf{notumor} & \textbf{pituitary} & \textbf{Total} \\
\midrule
Training (full) & 1,400 & 1,400 & 1,400 & 1,400 & 5,600 \\
Train (80\%) & 1,120 & 1,120 & 1,120 & 1,120 & 4,480 \\
Validation (20\%) & 280 & 280 & 280 & 280 & 1,120 \\
Test (held out) & 400 & 400 & 400 & 400 & 1,600 \\
Grand total & 1,800 & 1,800 & 1,800 & 1,800 & 7,200 \\
\bottomrule
\end{tabularx}
\end{center}
\section{3.6 Exploratory Data Analysis}
\begin{figure}[H]
  \centering
  \includegraphics[width=0.8\textwidth]{Figure_03_image_size_distribution.png}
  \caption{Distribution of original image dimensions.}
\end{figure}
\begin{figure}[H]
  \centering
  \includegraphics[width=0.8\textwidth]{Figure_05_pixel_intensity_histograms.png}
  \caption{Per-class pixel-intensity histograms showing substantial overlap.}
\end{figure}
The data set was first analysed based on four different criteria before training any models, so that any anomalies that would need to be dealt with could be flagged up.
File integrity: Each of the image files in the split was successfully opened using the Python Imaging Library. There were no problems with corrupt, incomplete, or invalid files.
Image formats: The dataset contains heterogeneity in terms of color format. From the total number of 7,200 images, 2,585 are stored in RGB format, 1,892 are in single channel greyscale format, and 3 are in four channel RGBA format. Conversion into RGB is performed at the beginning of each pipeline.
Size: The size of the image data varies from 150 x 168 pixels to 1,375 x 1,446 pixels. The average image has around 512 x 512 pixels. Images are standardized to the size of 224 x 224 pixels in preprocessing.
Distribution of pixel-intensity: There is substantial overlap between the histograms of pixel intensities of each class. In other words, pixel intensities alone cannot be used for separating any of the four classes of images.
\section{3.7 Train / Validation / Test Split Strategy}
A central methodological principle of this project is to divide the available data into three sets, the roles of which are strictly maintained during the project lifecycle:
Training set (4,480 images, 62.2\% of the corpus): Used for gradient changes during models that fit .
Validation set (1,120 images, 15.6\%): Implemented for early-stopping signals, hyper-parameter choice, ensemble-weight grid search and any other model choice decisions.
Test set (1,600 images, 22.2\%):  Served as fully until all the hyperparameters and modelling options are chosen, then analyzed exactly once to give the final results.
The train/validation split takes place by using sklearn.model\_selection.train\_test\_split with stratify=labels and random\_state=42, resulting in both class balance and bit-exact consistency. The resulting split is saved as csv-files in data/processed/ and is version-controlled alongside the original code.
The difference in validation-to-test accuracy throughout the experiments, which is between 4 and 8 \% throughout all three model types, indicates that any system presenting just validation-set metrics could significantly exaggerate its real-world performance.
\section{3.8 Image Preprocessing Pipeline}
\begin{figure}[H]
  \centering
  \includegraphics[width=0.8\textwidth]{Figure_06_preprocessing_effect.jpg}
  \caption{Visual effect of the deterministic preprocessing pipeline.}
\end{figure}
By applying a single deterministic preprocessing procedure to confirmation and test images, and to all images at the same time. The training process uses the same transform along with an additional stochastic augmentation phase described in Sec. 3.9.
Table 3.3: Image preprocessing transforms.
\begin{center}
\small
\begin{tabularx}{\textwidth}{YYY}
\toprule
\textbf{Step} & \textbf{Transform} & \textbf{Parameters} \\
\midrule
1 & Convert to RGB & PIL.Image.convert('RGB') \\
2 & Resize & 224 x 224 px, bilinear interpolation \\
3 & Convert to tensor & float32, range [0, 1], channels-first \\
4 & Normalise (ImageNet) & mean=[0.485, 0.456, 0.406], std=[0.229, 0.224, 0.225] \\
\bottomrule
\end{tabularx}
\end{center}
224 x 224 is selected because of the number of published ImageNet checkpoints for transfer learning. This resolution was used to pretrain both EfficientNet-B0 and Swin-Tiny, and a different target would destroy the pretrained positional embeddings (in the example of Swin) or the spatial requirements of early convolutional layers. In step 4 we use the ImageNet mean and std in order to make sure the inputs are in the distribution that the pretrained weights were trained on.
\section{3.9 Data Augmentation Strategy}
\begin{figure}[H]
  \centering
  \includegraphics[width=0.85\textwidth]{Figure_07_augmentation_examples.jpg}
  \caption{Examples of training-time data augmentation.}
\end{figure}
The augmentation is only used on the training images and only during training. The validation and the test images are never augmented. The augmentation policy is knowingly conservative, clearly not including any augmentation that results in clinically unrealistic images.
Table 3.4: Training-time augmentation policy.
\begin{center}
\small
\begin{tabularx}{\textwidth}{YYY}
\toprule
\textbf{Augmentation} & \textbf{Probability / Range} & \textbf{Clinical Justification} \\
\midrule
Random horizontal flip & p = 0.5 & Brain is bilaterally symmetric on axial slices \\
Random rotation & +/- 15 degrees & Simulates patient head tilt \\
Random affine (translate) & +/- 5\% & Simulates field-of-view variation \\
Random affine (scale) & +/- 5\% & Simulates magnification differences \\
Colour jitter (brightness) & +/- 0.20 & Simulates scanner-gain variation \\
Colour jitter (contrast) & +/- 0.20 & Simulates intensity-window variation \\
Random erasing & p=0.25, area 2-10\% & Forces model to reason about global structure \\
\bottomrule
\end{tabularx}
\end{center}
Three augmentations were looked at and carefully left out. Vertical flips might invert cranio-caudal anatomy. If you apply extreme rotations more than 30 degrees, the images will look unreal. In pilot experiments, we found that hue jitter reduces the validation accuracy, since the pixel intensities in MRI already have a diagnostic meaning.
\section{3.10 Ethical Considerations and Licensing}
The dataset is accessible to the public for educational and research use and does not contain any patient-identifiable information. No additional patient data were collected for this project. The result of the trained system is for academic demonstration purposes only and should not be offered as a medical device or a replacement for the judgment of a qualified radiologist. A long-term advisory to this effect is shown on the Streamlit interface. By default, the low-confidence threshold gate is set to suggest a radiologist review for any prediction with confidence below 85\%.
\chapter{Chapter 4: Project Planning}
\section{4.1 Overview}
Project planning is a blueprint of the creation and completion of the Brain Tumor Detection System. The planning process concentrates on separating the project into well-defined stages to ensure systematic progress, efficient management of time and clear deliverables for each phase.
\section{4.2 Current Status of the Project}
The project is now fully executed and functional at this point in time. The following tasks have been properly carried out:
Identification of problem area and creation of objectives of the project
Literature survey and analysis of existing systems
Detailed Dataset Acquisition Exploratory Data Analysis and Design of Preprocessing Pipeline
Planning and training of 3 deep learning models: Customized CNN, EfficientNet-B0 and Swin-Tiny
Ensemble designing with grid search weight optimization
A thorough assessment on the 1,600 images in the held-out test set
Grad-CAM explaining implementation for all 3 models
Creating Streamlit web interface and FastAPI REST backend
ONNX model export with inference verification 
Cloud deployment with Docker containerization
Full report writing and documentation
\section{4.3 System Architecture}
The Brain Tumor Detection system is constructed with a layered architecture to allow modularity, maintainability and ease of deployment. The system is composed of four logical layers:
Presentation Layer: This layer has two user-facing interfaces. The Streamlit web application provides a collaborative demo with drag-and-drop image upload, live forecasting view, confidence based metrics, per-model chance breakdowns, and Grad-CAM heatmap overlays. The FastAPI REST backend offers the same prediction functionality as a JSON API, which can be simple to integrate with hospital PACS systems, mobile apps or third-party scripts. 
Application Layer: The primary processing logic is in the src/ Python package. All prediction requests are produced via the single entry point, the class BrainTumorPredictor in src/inference.py. This loads all three trained models at once, performs each model on the input image, blends their outputs via the ensemble combiner and optionally produces Grad-CAM heatmaps. We have the same class received and called in the Streamlit app and the FastAPI backend so the actions are consistent across interfaces.
Model Layer The three trained model framework (src/cnn\_model.py, src/transfer\_model.py, src/swin\_model.py), an ensemble combiner (src/ensemble.py), the Grad-CAM execution (src/explainability.py), and the training engine (src/train.py). All model weights are stored as .pth checkpoint files within the models/ directory, along with an ensemble\_config.json that keeps track of the ensemble method and best weights.
Data Layer: A dataset can be found in DATASET/. The folders are Training/ and Testing/. data/processed/ contains cached divided split CSVs. The results of training are saved in outputs/ and include plots, confusion matrices, Grad-CAM images and CSV reports. Write TensorBoard event files to logs/. 
\section{4.4 Identified Challenges and System Limitations}
GPU memory limitation: The full training pipeline must be able to fit into 4 GB of VRAM (NVIDIA GTX 1650) . This keeps the largest batch size to 32 for Swin-Tiny in mixed precision.
Mixed-precision numerical instability: fp16 training caused a NaN bug in cuDNN on the Turing GPU architecture, leading to a switch to bfloat16.
Dominant-model effect: One baseline model (EfficientNet-B0) is much stronger than the other two, reducing the efficiency of ensembling.
Single MRI sequence: The system makes the categorization based on a single T1-weighted slice, while in clinical practice multiple MRI sequences are used for decision making.
\section{4.5 Strategy to Address Identified Challenges}
Bfloat16 mixed precision: Switching from fp16 to bfloat16 fixed NaN problem and kept the benefits of mixed-precision training.
Extremely aggressive but still clinically valid augmentation: Augmentation transforms are like multiplying your training data 7x, helping prevent overfitting on such a small dataset.
Use different architectures: Using 3 different architectures( vanilla CNN, depthwise - separable CNN, and transformer) helps ensure errors can be uncorrelated.
Grad-CAM support: Giving a heatmap explanation with each prediction helps address the clinical trust issue.
Threshold gate: Having a configurable threshold for prediction confidence allows uncertain predictions to be flagged for radiologist review instead of being output with wrong confidence. 
\section{4.6 Implementation Roadmap}
Table 4.1: Phases of the entire project work.
\begin{center}
\small
\begin{tabularx}{\textwidth}{YYY}
\toprule
\textbf{Phase} & \textbf{Module} & \textbf{Tasks} \\
\midrule
1 & Data Acquisition and EDA & Download dataset, verify integrity, compute class distributions, analyse image modes and dimensions, build stratified train/val/test split \\
2 & Custom CNN Training & Design 4-block VGG-style CNN, implement training engine with AMP and early stopping, train on 4,480 images, evaluate on validation set \\
3 & Transfer Learning (EfficientNet-B0) & Load ImageNet-pretrained backbone, replace classifier head, two-stage training (freeze then fine-tune), evaluate on validation set \\
4 & Swin Transformer Fine-tuning & Load timm Swin-Tiny checkpoint, replace head, two-stage training with very small Stage 2 LR, evaluate on validation set \\
5 & Ensemble Construction & Implement soft voting, weighted average, and stacking combiners, run exhaustive grid search for optimal weights, select best method \\
6 & Test-Set Evaluation & Run all three models and ensemble on held-out 1,600 test images, compute accuracy, F1, AUC, confusion matrices, ROC curves \\
7 & Explainability (Grad-CAM) & Implement Grad-CAM for all three architectures, generate heatmap overlays, analyse failure modes \\
8 & Deployment & Build Streamlit UI, FastAPI REST endpoint, ONNX export pipeline, Docker container \\
9 & Documentation & Write project report, prepare viva QA guide, generate flow diagrams and architecture documentation \\
\bottomrule
\end{tabularx}
\end{center}
\section{4.7 Project Scheduling and Gantt Chart}
\begin{figure}[H]
  \centering
  \includegraphics[width=0.92\textwidth]{Figure_31_gantt_chart.png}
  \caption{Project schedule across the semester.}
\end{figure}
The project took place during one semester. Figure 4.2 shows the Gantt chart detailing schedule and dependencies between tasks. Steps 2-4 (model training) were concurrent to a degree - the custom CNN was trained first as a baseline, then the two transfer-learning models. Ensemble and evaluation could not be parallelised since they require all three models to be finished training.
\chapter*{}
\addcontentsline{toc}{chapter}{}
\chapter{Chapter 5: Project Description}
\section{5.1 Model of the Software}
The Brain Tumour Detection system is developed using modular python package structure. All production code lives in src/ and is broken into modules with a single responsibility. The two interfaces that users will see are in the app/ folder. The notebooks/ directory contains seven Jupyter notebooks that reproduce the entire experimental pipeline with embedded outputs for auditability.
Here is a summary of the key technology decisions and rationale:
\begin{center}
\small
\begin{tabularx}{\textwidth}{YYY}
\toprule
\textbf{Component} & \textbf{Choice} & \textbf{Reason} \\
\midrule
Framework & PyTorch 2.12 & Industry standard for research; native CUDA support \\
Mixed precision & bfloat16 & Avoids fp16 cuDNN NaN bug on Turing GPUs \\
Optimiser & AdamW & Strong baseline for small medical datasets \\
Loss function & CrossEntropyLoss & Classes are balanced; no weighting needed \\
Batch size & 32 & Largest that fits Swin-Tiny in 4 GB VRAM with AMP \\
LR scheduler & CosineAnnealingLR & Smooth convergence with minimal tuning \\
Random seed & 42 (all RNGs) & Bit-exact reproducibility \\
Hardware & NVIDIA GTX 1650 (4 GB), Windows 10, Python 3.11 &  \\
\bottomrule
\end{tabularx}
\end{center}
\section{5.2 Software Requirements Specification (SRS)}
5.2.1 Introduction
Brain Tumour Detection is a deep learning based diagnostic aid which classifies brain MRI slices into 4 categories. This SRS document describes the functional and non functional requirements, system constraints and expected performance characteristics
5.2.2 General description
The system is a client-server application with a web-based front-end and a backend written in python. It has two modes of interaction,
The system is a client-server application with a backend written in python and a web-based frontend. It has two modes of interaction,
(1) Interactive single image prediction with visual feedback using Streamlit web interface
(2) Programmatic batch predictions with a FastAPI REST endpoint. forecast. Both interfaces wrap the same BrainTumorPredictor class, so accessing the class through either interface will give the same results.
\subsection{5.2.3 Functional Requirements}
The system will be:
Output: Brain MRI Image in JPG, PNG or BMP Image format. The original image size is not limited.
Convert the input image to RGB, resize to 224 x 224 and normalise using ImageNet statistics.
Run the preprocessed image through 3 trained models (custom CNN, EfficientNet-B0, Swin-Tiny). Output softmax probability vectors for each.
Then the three probability vectors are combined to a single ensemble probability distribution using the saved ensemble method and weights.
Return predicted class (class with max ensemble probability) and confidence score of this class
(Optional) Create Grad-CAM heatmap overlays for the predicted class for each base model.
Flag predictions with less than a user configurable threshold of confidence (default 0.85) with "consult radiologist" advice
\begin{itemize}[leftmargin=*]
  \item Expose a /predict REST endpoint that accepts multipart/form-data image upload and returns a JSON object with predicted class, confidence and per-model probabilities. Streamlit UI: predicted class (colour coded), confidence, probability bar chart, probability table per model, 3 Grad-CAM heatmap images.
\end{itemize}
\subsection{5.2.4 Interface Requirements}
User Interface:
Web-based interface that runs on any modern browser.
Drag \& Drop image upload or choose test-sample from dropdown.
Model info \& settings in sidebar.
Hardware Interface:
Runs on a regular desktop or laptop, with or without dedicated GPU. Back to CPU gracefully.
Software Interface:
Python 3.11+, PyTorch 2.12, timm, Streamlit, FastAPI, ONNX Runtime
\subsection{5.2.5 Performance Requirements}
Single image inference should be performed in 300 ms on GPU (including Grad-CAM) or 1 second on CPU.
The first startup should take less than 10 seconds to load the model.
The Streamlit interface should remain responsive during prediction.
\subsection{5.2.6 Design Constraints}
The system will be developed using open source frameworks and tools.
All the model weights are going to be stored locally, no external API calls needed at inference time.
The system shall be deployable on standard institutional infrastructure or personal laptops.
\subsection{5.2.7 Non-Functional Requirements}
Reliability: The system should be deterministic and reliable for the same input image. For reproducibility we will use static random seeds and cache data splits.
Security: No user data will be stored or sent anywhere else. Image inputs are only stored in memory, not persistently.
Maintainability: Modularity \& modular src/ folder structure to allow easy modifications. Retraining and substituting each model does not affect other pieces of code.
Scalability : FastAPI framework allows concurrent connection. Exporting to Onnx will allow models to be run on edge devices without the PyTorch framework.
\section{5.3 Functional Specification}
\subsection{5.3.1 Deep Learning Models}
\begin{figure}[H]
  \centering
  \includegraphics[width=0.92\textwidth]{Figure_09_architecture_custom_cnn.png}
  \caption{Custom CNN architecture.}
\end{figure}
\begin{figure}[H]
  \centering
  \includegraphics[width=0.92\textwidth]{Figure_10_architecture_efficientnet_b0.png}
  \caption{EfficientNet-B0 transfer-learning architecture.}
\end{figure}
\begin{figure}[H]
  \centering
  \includegraphics[width=0.92\textwidth]{Figure_11_architecture_swin_tiny.png}
  \caption{Swin-Tiny architecture.}
\end{figure}
Model 1: Custom CNN (src/cnn\_model.py)
VGG style architecture, 4 convolutional blocks. We have two consecutive 3x3 convolutions in each block. After each convolution, a Batch Normalisation and ReLU activation is applied. The 2x2 max-pooling layer reduces the spatial dimensions by half. The number of channels increases by a factor of 2 at each block: 32, 64, 128, 256. The feature map (256 x 14 x 14) generated by the fourth block is fed into a Global Average Pooling layer, which reduces the spatial dimensions and outputs a 256 dimensional feature vector. This vector is fed to a two-layer fully-connected classifier: Dropout Linear (256 to 128) ReLU Dropout Linear (128 to 4) which outputs four logits. Weight initialisation Kaiming normal (mode=fan\_out, relu) Total parameters: 1,206,628
Model 2: EfficientNet-B0 Transfer Learning (src/transfer\_model.py)
We chose EfficientNet-B0 over ResNet-50, EfficientNet-B3, DenseNet-121 and ConvNeXt-Tiny for its accuracy-per-parameter efficiency (77.7\% ImageNet top-1 at only 5.3 million parameters) and its suitability for the 4 GB VRAM budget. The pretrained backbone comes from torchvision. We remove its 1000-class classifier head and replace it with a new one: Dropout(0.3) -> Linear(1280 to 4).
Training is done in two steps. Stage 1 (warm-up, 4 epochs, learning rate = 0.001): The backbone is frozen and only the 5,124 head parameters are updated with gradients. We override the module's .train() method to keep the frozen backbone in .eval() mode, to avoid BatchNorm running-statistic drifts. Unfreeze backbone Stage 2 (fine-tune, 8 epochs, learning rate = 0.0001) The 10x smaller learning rate helps to mitigate the catastrophic forgetting of pre-trained features. Total params: 4,012,672
Model 3: Swin-Tiny Transformer (src/swin\_model.py)
Swin-Tiny (swin\_tiny\_patch4\_window7\_224 from the timm library) is a hierarchical vision transformer that accepts an image composed of 4 x 4 patches as input and then processes it through four stages of windowed self-attention blocks. In each stage , spatial resolution decreases by merging patches ( like pooling in CNNs ) , while feature dimensionality increases . Swin's key innovation is the shifted-window mechanism, where attention windows are shifted by half a window size in alternating layers, allowing information to flow between windows without the quadratic costs of global attention.
The classifier head is replaced as in EfficientNet (Dropout 0.3, Linear 768 to 4). Two-stage training follows the transfer-learning recipe but with a much smaller Stage 2 learning rate. rate. of 0.00002 (50 times smaller than Stage 1) because transformer attention weights are very sensitive to large gradient updates. Total parameters: 27,522,430.
Table 5.3: Architectural summary.
\begin{center}
\small
\begin{tabularx}{\textwidth}{YYYY}
\toprule
\textbf{Property} & \textbf{Custom CNN} & \textbf{EfficientNet-B0} & \textbf{Swin-Tiny} \\
\midrule
Total params & 1.21 M & 4.01 M & 27.52 M \\
Architecture type & VGG-style conv blocks & Depthwise-separable + SE & Shifted-window self-attention \\
Pretrained & No (from scratch) & Yes (ImageNet) & Yes (ImageNet) \\
Stage 1 LR & 0.001 & 0.001 & 0.001 \\
Stage 2 LR & N/A (single stage) & 0.0001 & 0.00002 \\
Stage 1 epochs & 20 & 4 & 3 \\
Stage 2 epochs & N/A & 8 & 4 \\
\bottomrule
\end{tabularx}
\end{center}
\subsection{5.3.2 Mathematical Model and Loss Function}
All three models are trained with standard cross-entropy loss. For a single sample with true class y and predicted probability p\_y for that class, the loss is:
    L=-log(p\_y)
where p\_y is the probability of the correct class output by softmax. The loss is not class-weighted, since the four classes are perfectly balanced.
The optimiser used for all three models is AdamW (Adam with decoupled weight decay). The main hyperparameters are: beta\_1 = 0.9, beta\_2 = 0.999, and weight decay = 1e-4. The learning rate schedule is CosineAnnealingLR which reduces the learning rate gradually from the initial value to eta\_min = initial\_lr x 0.01 during training. Mixed precision training uses bfloat16 through torch.amp.autocast which reduces memory by 2x and increases throughput on compatible GPUs without the numerical instability issues of fp16 on Turing-architecture GPUs.
\subsection{5.3.3 Ensemble Decision Logic}
\begin{figure}[H]
  \centering
  \includegraphics[width=0.7\textwidth]{Figure_27_ensemble_weight_heatmap.png}
  \caption{Validation accuracy across the ensemble-weight grid search.}
\end{figure}
Each trained model outputs a softmax probability vector p\_i  in  R\textasciicircum{}4 (one probability for each class). Three ensemble combination methods were tested:
Soft voting: The ensemble probability is the average of the three model outputs, p\_ens = (p\_cnn + p\_tl + p\_swin)/3.
Weighted average: The ensemble probability is a weighted sum: p\_ens = w\_cnn x p\_cnn + w\_tl x p\_tl + w\_swin x p\_swin with the weights w\_i >= 0 and w\_cnn + w\_tl + w\_swin = 1. The optimal weights were found by exhaustive grid search over the simplex at step size 0.05, yielding 231 candidate weight combinations for three models. Each combination was evaluated on the validation set, and the best-performing combination was selected: [w\_cnn=0.0, w\_tl=0.9, w\_swin=0.1].
Stacking: A logistic regression meta-learner was trained on the concatenated 12-dimensional probability vector (4 probabilities from each of 3 models) using 5-fold cross-validation on the validation set.
The grid search effectively discarded the custom CNN (w\_cnn=0.0) because its predictions overlap too much with EfficientNet's. Adding the CNN brings no complementary signal but does add noise. The weighted average method was selected as the final ensemble combiner based on its validation-set performance.
\subsection{5.3.4 Grad-CAM Explainability Generation}
\begin{figure}[H]
  \centering
  \includegraphics[width=0.9\textwidth]{Figure_28_gradcam_correct.jpg}
  \caption{Grad-CAM overlays for correctly classified samples (one per class).}
\end{figure}
Gradient-weighted Class Activation Mapping (Grad-CAM) is an approach to generate a coarse localisation heatmap highlighting the important regions in the input image for predicting the concept. It operates by calculating the gradient of the scores for the target class w.r.t the feature maps of a selected convolutional layer, and then combining the feature maps linearly using these gradients [13].
For the custom CNN and EfficientNet-B0, the target layer is the last convolutional layer before global average pooling. The target is the LayerNorm at the end of the last transformer stage for Swin-Tiny without the traditional convolutional layers. Swin's activations are in channels-last format (sequence\_length x channels) while Grad-CAM expects them in channels-first format (channels x height x width), so a reshape is necessary. The implementation is based on pytorch-grad-cam library by Jacob Gildenblat.
The heatmap produced is a single channel gray-scale image, with red indicating high importance and blue indicating low importance. This heatmap is then superimposed on the original MRI image at 45\% opacity to create a composite visualisation that allows the clinician to simultaneously view the anatomical structure and the attention pattern of the model.
\section{5.4 Design Specification}
\subsection{5.4.1 Use Case Diagrams}
\begin{figure}[H]
  \centering
  \includegraphics[width=0.8\textwidth]{Figure_15_usecase_radiologist.png}
  \caption{Use-case diagram: Radiologist / Clinician role.}
\end{figure}
\begin{figure}[H]
  \centering
  \includegraphics[width=0.8\textwidth]{Figure_16_usecase_researcher.png}
  \caption{Use-case diagram: Researcher / Developer role.}
\end{figure}
\begin{figure}[H]
  \centering
  \includegraphics[width=0.8\textwidth]{Figure_17_usecase_admin.png}
  \caption{Use-case diagram: Administrator role.}
\end{figure}
The system supports three main actor roles:
Radiologist / Clinician: Upload a patient MRI Image (or choose test sample). View the predicted class and its confidence. Check the Grad-CAM heatmaps to make sure that the model is reasoning correctly. Adjust the confidence threshold.
Researcher / Developer Calls batch prediction API of FastAPI Looks at per-model probability breakdowns Runs evaluation scripts on custom datasets Exports models to ONNX
Administrator: Launches the docker container, configures the server, monitors the application logs, and updates the model checkpoints.
\subsection{5.4.2 Data flow diagrams}
\begin{figure}[H]
  \centering
  \includegraphics[width=0.92\textwidth]{Figure_18_dfd_training.png}
  \caption{Data Flow Diagram: model training pipeline.}
\end{figure}
\begin{figure}[H]
  \centering
  \includegraphics[width=0.92\textwidth]{Figure_30_dfd_inference.png}
  \caption{Data Flow Diagram: single-image inference.}
\end{figure}
DFD of training pipeline: Data loader (src/data\_loader.py) takes raw MRI images from DATASET/ and applies preprocessing and augmentation transforms (src/preprocess.py). Then the processed batches are fed to the training engine (src/train.py) which does forward passes, computes the loss, backpropogates and logs metrics to tensorboard. Engine does evaluation on validation set in every epoch and saves the best checkpoint to `models/`
Inference Pipeline DFD User uploads an image through Streamlit UI or FastAPI endpoint. The image is passed to BrainTumorPredictor ( src/inference.py) which loads the image, applies the inference transform, runs all three models in parallel, combines the softmax outputs with the ensemble combiner and optionally produces Grad-CAM overlays. The prediction result object is passed on to the presentation layer for visualisation.
\subsection{5.4.3 Data Dictionaries}
The following tables describe the main data structures used in the system.
Table 5.6: Ensemble Configuration Schema (ensemble\_config.json).
\begin{center}
\small
\begin{tabularx}{\textwidth}{YYY}
\toprule
\textbf{Field} & \textbf{Type} & \textbf{Description} \\
\midrule
method & string & Ensemble method: 'weighted', 'soft\_voting', or 'stacking' \\
weights & list of float & Mixing weights for the three models, summing to 1.0 \\
val\_accuracy & float & Validation accuracy achieved with these weights \\
models & list of string & Ordered list of model names: ['cnn', 'transfer', 'swin'] \\
\bottomrule
\end{tabularx}
\end{center}
Table 5.7: Prediction Result Schema.
\begin{center}
\small
\begin{tabularx}{\textwidth}{YYY}
\toprule
\textbf{Field} & \textbf{Type} & \textbf{Description} \\
\midrule
predicted\_class & string & One of: glioma, meningioma, pituitary, notumor \\
predicted\_class\_idx & int & Integer index (0-3) \\
confidence & float & Ensemble top-class probability \\
ensemble\_probs & dict & Class name to probability mapping \\
model\_probs & dict of dict & Per-model, per-class probs \\
inference\_time\_ms & float & Wall-clock inference time \\
gradcams & dict or None & Model name to heatmap array (if requested) \\
\bottomrule
\end{tabularx}
\end{center}
Table 5.8: Model Checkpoint Registry.
\begin{center}
\small
\begin{tabularx}{\textwidth}{YYY}
\toprule
\textbf{File} & \textbf{Size} & \textbf{Contents} \\
\midrule
cnn\_model.pth & 4.9 MB & Custom CNN state\_dict \\
transfer\_model.pth & 16.4 MB & EfficientNet-B0 state\_dict \\
swin\_model.pth & 110.2 MB & Swin-Tiny state\_dict \\
cnn\_model.onnx & 4.8 MB & ONNX export of custom CNN \\
transfer\_model.onnx & 16.0 MB & ONNX export of EfficientNet-B0 \\
swin\_model.onnx & 112.9 MB & ONNX export of Swin-Tiny \\
ensemble\_config.json & 0.4 KB & Ensemble method and weights \\
\bottomrule
\end{tabularx}
\end{center}
Table 5.11: Inference API Schema (/predict).
\begin{center}
\small
\begin{tabularx}{\textwidth}{YYYY}
\toprule
\textbf{Direction} & \textbf{Field} & \textbf{Type} & \textbf{Description} \\
\midrule
Request & file & multipart/form-data & MRI image file \\
Response & predicted\_class & string & Predicted category \\
Response & confidence & float & Top-class probability \\
Response & ensemble\_probs & object & All 4 class probs \\
Response & model\_probs & object & Per-model breakdown \\
Response & inference\_time\_ms & float & Processing time \\
\bottomrule
\end{tabularx}
\end{center}
\chapter{Chapter 6: Implementation Issues}
\section{6.1 Diversity of Datasets and Image Quality}
The first problem that appeared during the implementation was the variety of source images. The Nickparvar collection consists of three different data sets, indicating different colour mode (RGB, greyscale, RGBA), dimensions (150 x 168 to 1,375 x 1,446 pixels), illumination intensity and contrast ratios of the pictures . Such variety is good from the perspective of future model performance, but this variety had to be processed carefully so that each image entered the model in the same format (3-channel, 224 × 224, ImageNet-normalised).
\section{6.2 Model choice and hyperparameter sensitivity}
\begin{figure}[H]
  \centering
  \includegraphics[width=0.8\textwidth]{Figure_12_cnn_training_curves.png}
  \caption{Custom CNN training and validation curves.}
\end{figure}
\begin{figure}[H]
  \centering
  \includegraphics[width=0.8\textwidth]{Figure_13_efficientnet_training_curves.png}
  \caption{EfficientNet-B0 training curves across both stages.}
\end{figure}
\begin{figure}[H]
  \centering
  \includegraphics[width=0.8\textwidth]{Figure_14_swin_training_curves.png}
  \caption{Swin-Tiny training curves across both stages.}
\end{figure}
In the first round of backbone selection, five candidates were selected for transfer-learning based backbones, including ResNet-50, EfficientNet-B0, EfficientNet-B3, DenseNet-121 and ConvNeXt-Tiny. We chose EfficientNet-B0 for its efficiency in accuracy per parameter and that it could fit under the 4 GB VRAM capacity limit. EfficientNet-B3 was slightly better in the validation accuracy but could not fit in the memory when the batch size is 32, so we had to reduce it to batch size 16 which resulted in a longer training time.
The Swin Transformer was quite sensitive to the learning rate in Stage 2 (Fine-tuning). A Learning Rate of 0.0001 in stage 2 (used for EfficientNet) resulted in collapse of validation accuracy in just 3 epochs. I was able to solve this with a much smaller Learning Rate for Stage 2 - 0.00002 (50x lower than stage 1). This sensitivity arises from the fact that the self-attention weights of transformers encode global relationships that are susceptible to large gradient updates.
\section{6.3 GPU Memory Limits and Mixed-Precision}
The whole training pipeline had to run on 4 GB VRAM of NVIDIA GTX 1650 (Turing arch, compute capability 7.5). Initial implementation used mixed precision training via fp16 to reduce VRAM usage and increase throughout. However, when training Swin Tiny, the network output NaN loss after \textasciitilde{}5 epochs.
By layer by layer activation tracing, we traced the root cause of the problem and found that it was due to numerical instability of cuDNN LayerNorm operations on Swin Transformer in fp16 on Turing arch. Turns out the solution to the problem was to use bfloat16 mixed precision instead of fp16. Bfloat16 has 16 bits , same as fp16 , but it has 8 exponent bits , same as float32 , so it has the same dynamic range .
.
\section{6.4 Class Confusion and Generalisation Gap}
\begin{figure}[H]
  \centering
  \includegraphics[width=0.6\textwidth]{Figure_19_confusion_matrix_cnn.png}
  \caption{Test-set confusion matrix: Custom CNN.}
\end{figure}
\begin{figure}[H]
  \centering
  \includegraphics[width=0.6\textwidth]{Figure_20_confusion_matrix_efficientnet.png}
  \caption{Test-set confusion matrix: EfficientNet-B0.}
\end{figure}
\begin{figure}[H]
  \centering
  \includegraphics[width=0.6\textwidth]{Figure_21_confusion_matrix_swin.png}
  \caption{Test-set confusion matrix: Swin-Tiny.}
\end{figure}
\begin{figure}[H]
  \centering
  \includegraphics[width=0.6\textwidth]{Figure_22_confusion_matrix_ensemble.png}
  \caption{Test-set confusion matrix: Weighted ensemble.}
\end{figure}
\begin{figure}[H]
  \centering
  \includegraphics[width=0.8\textwidth]{Figure_24_per_class_f1_bars.png}
  \caption{Per-class F1 scores across all models.}
\end{figure}
From the confusion matrices analysis, it seems that the main cause for most of the classification errors for the three networks is the confusion between glioma and meningioma. The custom CNN confuses gliomas with meningiomas in 17\% of the cases in contrast to 1\% and 6\% with EfficientNet-B0 and Swin-Tiny, respectively. Clinically, the two tumours may appear almost identical on a single T1-weighted axial image, which makes their confusion plausible.
The gap between the three networks in terms of the generalisation gap between validation set and test set accuracy ranges from 4\% to 8\%. This generalisation gap is a known fact in small-sized datasets in medical imaging and it emphasises the importance of reporting test-set accuracy rather than validation-set accuracy. Such a system would report only validation accuracy and would be misleading.
\section{6.5 Trustworthiness in Clinical Adoption and Explainability}
\begin{figure}[H]
  \centering
  \includegraphics[width=0.9\textwidth]{Figure_29_gradcam_misclassified.jpg}
  \caption{Grad-CAM overlays for misclassified samples, illustrating failure modes.}
\end{figure}
The major hurdle in the adoption of AI based diagnostic systems is not accuracy, but trustworthiness, in terms of practical implementation. For clinicians to use AI diagnostics, they need to be sure of the reason why the model made the specific decision it did. This is achieved by the implementation of three components: (1) the Grad-CAM visualisation of MRI hot spots analysed by the algorithm, (2) comparison of individual probability estimates from each of the three models to see if they coincide, (3) adjustable probability thresholds allowing only high-confidence predictions to be used for diagnosis.
Analysis of Grad-CAM heat maps of correctly classified examples demonstrates that the models attend to anatomically relevant locations (tumour area for tumour classification and brain parenchyma for no-tumor cases). Heat maps of EfficientNet show the highest focus and Swin the least focused, in line with the larger effective receptive fields of self-attention. The heatmap analysis of the misclassified samples reveals the model correctly identifies the anatomical location, but makes a wrong interpretation of the visual features, which indicates clinical ambiguity and not misleading features.
\section{6.6 Problem of Dominant Model in Ensemble}
\begin{figure}[H]
  \centering
  \includegraphics[width=0.75\textwidth]{Figure_23_roc_curves_test.png}
  \caption{ROC curves on the test set for all models and the ensemble.}
\end{figure}
\begin{figure}[H]
  \centering
  \includegraphics[width=0.8\textwidth]{Figure_25_final_comparison_bars.png}
  \caption{Final comparison of test-set accuracy, F1 and macro AUC.}
\end{figure}
\begin{figure}[H]
  \centering
  \includegraphics[width=0.75\textwidth]{Figure_26_efficiency_pareto.png}
  \caption{Accuracy versus inference-latency trade-off.}
\end{figure}
The most important and unexpected empirical result of this study is that the three-model ensemble (94.69\% test accuracy) does not outperform the best single model (EfficientNet-B0, 94.94\% test accuracy). The result is opposite to the usual expectation for an ensemble, namely, the improvement over each of its components. As such, a closer examination of the situation seems reasonable.
There are two reasons for the failure mentioned above. First of all, the efficientNet-B0 works very close to the noise ceiling of the dataset, i.e., has a validation error of 1.07\% and a test error of 5.06\%. Thus, to surpass it, complementary errors from the base models would be required. Second, the errors made by the custom CNN and swin-tiny are quite similar to those of the efficientNet-B0 (especially on the glioma-meningioma boundary), meaning that their prediction averaging does not reduce errors but introduces extra noise due to weak models.
The phenomenon described above was named 'the dominant-model problem' and discussed in the field of ensemble learning. Namely, the variance reduction promised by the ensemble technique only becomes possible if the errors are uncorrelated, which does not happen in our case, as all three models experience problems on the same boundary.
The reason we choose to report our finding is because it will have some practical implications, and also it will be good for pedagogical reasons. In terms of its practical applications, it shows us that in the process of constructing an ensemble, not only the architecture diversity is important, but also the error diversity should be taken into account.
Table 6.1: Per-model test-set performance.
\begin{center}
\footnotesize
\begin{tabularx}{\textwidth}{YYYYYYYY}
\toprule
\textbf{Model} & \textbf{Params} & \textbf{Val acc} & \textbf{Test acc} & \textbf{Macro F1} & \textbf{Macro AUC} & \textbf{Train time} & \textbf{Latency} \\
\midrule
Custom CNN & 1.21 M & 0.9036 & 0.8275 & 0.8221 & 0.9513 & 32.1 min & 6.8 ms \\
EfficientNet-B0 & 4.01 M & 0.9893 & 0.9494 & 0.9482 & 0.9908 & 14.9 min & 24.1 ms \\
Swin-Tiny & 27.52 M & 0.9670 & 0.9194 & 0.9171 & 0.9865 & 12.5 min & 25.7 ms \\
Ensemble (weighted) & 32.74 M & 0.9893 & 0.9469 & 0.9456 & 0.9907 & 59.5 min & 56.6 ms \\
\bottomrule
\end{tabularx}
\end{center}
Table 6.2: Ensemble grid-search results (top 5).
\begin{center}
\footnotesize
\begin{tabularx}{\textwidth}{YYYYY}
\toprule
\textbf{Rank} & \textbf{w\_cnn} & \textbf{w\_transfer} & \textbf{w\_swin} & \textbf{Val accuracy} \\
\midrule
1 & 0.00 & 0.90 & 0.10 & 0.9893 \\
2 & 0.00 & 0.95 & 0.05 & 0.9893 \\
3 & 0.00 & 1.00 & 0.00 & 0.9893 \\
4 & 0.05 & 0.90 & 0.05 & 0.9884 \\
5 & 0.00 & 0.85 & 0.15 & 0.9875 \\
\bottomrule
\end{tabularx}
\end{center}
\chapter{Chapter 7: Conclusion, Summary and Future Scope}
\section{7.1 Summary of Work}
This project developed a complete deep learning pipeline for the classification of brain tumours into four classes using T1-weighted MRI scans. This project has three different models with architectural diversity, a custom CNN model trained from scratch (1.2 million parameters), an EfficientNet-B0 model with two-stage transfer learning from ImageNet (4.0 million parameters), and a Swin-Tiny vision transformer model fine-tuned with appropriate learning rates (27.5 million parameters). Moreover, the three models are combined together by an optimal weighted-average ensemble with mixing weights [0.0, 0.9, 0.1] obtained after grid search on the validation data.
These three models are trained and tested on Brain Tumour MRI Dataset (Nickparvar, Kaggle) which has 7,200 images in four different categories. Stratification is used to divide the dataset into a training set with 4,480 images and a validation set with 1,120 images. The remaining 1,600 images are reserved for testing and are used only once to assess the metrics of interest.
The project has also provided deployment interfaces, Grad-CAM explanations, low-confidence gate and project documentation in addition to model development and evaluation.
\section{7.2 Conclusion}
The leading individual model, EfficientNet-B0 with transfer learning, achieves 94.94\% accuracy on the test data and a macro-averaged ROC-AUC value of 0.991. The three-model ensemble model reaches 94.69\% accuracy and 0.991 AUC, equal to the best performing individual model, but does not outperform it. These results are in agreement with the dominant model problem identified in the ensemble learning literature in that when there is a pronounced disparity between the performances of the learners such that some make correlated errors, ensembling introduces noise.
As observed in the Grad-CAM results, all three models attend to anatomically relevant areas. The heatmaps of EfficientNet-B0 have the most precise localisation, while those of Swin are the most scattered. The few remaining errors are in clinically ambiguous areas. In particular, the border between gliomas and meningiomas.
Deployability of the system is achieved via four distinct interfaces: Streamlit for interactive demo, FastAPI for programmatic access, ONNX model export for faster performance (2 to 3.5 times faster than plain CPU inference), and Docker for cloud-based deployment. All four interfaces utilize the same checkpoint models, thus providing equivalent prediction outcomes irrespective of how the system is accessed.
The project proves that an entire medical imaging AI system can be developed from scratch, trained, and deployed on relatively small equipment such as one NVIDIA GTX 1650 with 4 GB of VRAM in just one academic semester with appropriate focus on basic engineering principles: reproducible data splits, mixed precision training, two-step transfer learning, honest test set assessment, and visual explanation.
\section{7.3 Future Scope}
While the current system delivers a robust solution, several avenues for future enhancement exist:
3D Volumetric Classification: Presently, the approach works on each two-dimensional slice separately. Future research may include the development of the model in processing three-dimensional MRI volumes, employing approaches like 3D-ResNet or V-Net, in order to capture the context between slices and enhance tumor classification accuracy.
Multi-Sequence Fusion: Clinical radiology utilizes various MR sequence imaging techniques such as T1, T2, FLAIR, and contrast-enhanced T1 for diagnosis purposes. Involving several sequences as input channels or adopting a multi-stream fusion approach can greatly help in distinguishing between similar-looking tumors through MR sequence imaging alone..
Self-Supervised Pretraining: Training a foundational model on a huge dataset of unlabeled brain MRI images by employing masked image modeling, among other methods, before training with supervision can result in more domain-relevant feature representations compared to those trained on ImageNet.
Test-Time Augmentation (TTA):Inference latency could be traded off for better accuracy by averaging the outputs from the model applied to various augmentations of the same test image (such as horizontal flip and slight rotations).
Cross-Institution Validation: This current assessment is done on one open dataset only. Before going into actual clinical use, validation of this system using multiple clinical datasets from various scanning machines will be absolutely necessary.
Federated Learning:For overcoming the privacy issue in medical settings, the model training pipeline can be altered such that the models are trained using federated learning at hospitals rather than centralized patient information.
LLM-Assisted Report Generation: Integrating a large language model to generate preliminary radiology reports from the classification output and Grad-CAM findings could further assist radiologists in high-volume settings.
\chapter*{Bibliography}
\addcontentsline{toc}{chapter}{Bibliography}
[1] Bouhafra, S. and El Bahi, H., "Deep Learning Approaches for Brain Tumor Detection and Classification Using MRI Images (2020 to 2024): A Systematic Review," Journal of Digital Imaging Informatics in Medicine, vol. 38, pp. 1403-1433, 2025.
[2] Ranjbarzadeh, R. et al., "Explainable AI and Vision Transformers for Detection and Classification of Brain Tumor: A Comprehensive Survey," Artificial Intelligence Review, Springer, 2025.
[3] Babu Vimala, B., Srinivasan, S., Mathivanan, S.K. et al., "Detection and Classification of Brain Tumor Using Hybrid Deep Learning Models," Scientific Reports, vol. 13, 23029, 2023.
[4] Islam, M.T. et al., "BrainNet: Precision Brain Tumor Classification with Optimized EfficientNet Architecture," International Journal of Intelligent Systems, Wiley, 2024.
[5] Alnowami, M. et al., "Enhancing EfficientNetv2 with Global and Efficient Channel Attention Mechanisms for Accurate MRI-Based Brain Tumor Classification," Cluster Computing, Springer, 2024.
[6] Haq, A.U. et al., "A Novel Swin Transformer Approach Utilizing Residual Multi-Layer Perceptron for Diagnosing Brain Tumors in MRI Images," International Journal of Machine Learning and Cybernetics, vol. 15, pp. 3579-3597, 2024.
[7] Alsubai, S. et al., "Enhanced Magnetic Resonance Imaging-Based Brain Tumor Classification with a Hybrid Swin Transformer and ResNet50V2 Model," Applied Sciences, vol. 14, no. 22, 10154, 2024.
[8] Khan, M.A. et al., "Weighted Average Ensemble Deep Learning Model for Stratification of Brain Tumor in MRI Images," Diagnostics, vol. 13, no. 7, 1320, 2023.
[9] Abdella, G.M. et al., "Majority Voting Ensemble of Deep CNNs for Robust MRI-Based Brain Tumor Classification," Diagnostics, vol. 15, no. 14, 1782, 2025.
[10] Ali, M. et al., "Enhancing Brain Tumor Detection in MRI Images through Explainable AI Using Grad-CAM with ResNet 50," BMC Medical Imaging, 2024.
[11] Naser, M.A. and Deen, M.J., "Explainable Deep Learning Approach for Multi-Class Brain MRI Tumor Classification and Localization Using Gradient-Weighted Class Activation Mapping," Information, vol. 14, no. 12, 642, 2023.
[12] Nickparvar, M., "Brain Tumor MRI Dataset," Kaggle, 2022. Available: https://www.kaggle.com/datasets/masoudnickparvar/brain-tumor-mri-dataset
[13] Selvaraju, R.R. et al., "Grad-CAM: Visual Explanations from Deep Networks via Gradient-based Localization," in Proceedings of the IEEE International Conference on Computer Vision (ICCV), 2017.
[14] Tan, M. and Le, Q.V., "EfficientNet: Rethinking Model Scaling for Convolutional Neural Networks," in Proceedings of the International Conference on Machine Learning (ICML), 2019.
[15] Liu, Z. et al., "Swin Transformer: Hierarchical Vision Transformer Using Shifted Windows," in Proceedings of the IEEE International Conference on Computer Vision (ICCV), 2021.
[16] Loshchilov, I. and Hutter, F., "Decoupled Weight Decay Regularization," in Proceedings of the International Conference on Learning Representations (ICLR), 2019.
[17] International Association of Cancer Registries (IARC) and Indian Journal of Neurology, 2025; BW Healthcare World, "Battling the Brain Tumor Burden: Challenges and Advancements in India," 2025.
[18] Micikevicius, P. et al., "Mixed Precision Training," in Proceedings of the International Conference on Learning Representations (ICLR), 2018.
[19] Ahmed, S. et al., "Advancing Brain Tumor MRI Classification using SwRD: A Parallel Swin Transformer-ResNet Approach," Virtual Reality and Intelligent Hardware, vol. 7, no. 5, pp. 501-522, 2025.
[20] Global Cancer Observatory (GLOBOCAN), "Cancer Today: Estimated Number of New Cases in 2020, Worldwide," International Agency for Research on Cancer, World Health Organization, 2022.
\end{document}
